\documentclass{article}
\usepackage{pgfplots} % plots
\usepackage{listings} % code
\usepackage{geometry}
\usepackage{array}
\usepackage{booktabs}
\usepackage{multirow}

\geometry{a4paper, portrait, margin=1in}

\definecolor{comments}{rgb}{0,0.6,0}
\definecolor{keywords}{rgb}{0.70, 0, 0.85}
\definecolor{numbers}{rgb}{0.45,0.35,0.15}
\definecolor{strings}{rgb}{0.65,0.1,0.1}
\definecolor{background}{rgb}{0.95,0.95,0.95}

\lstdefinestyle{python_style}{
    backgroundcolor=\color{background},   
    commentstyle=\color{comments},
    keywordstyle=\color{keywords},
    numberstyle=\tiny\color{numbers},
    stringstyle=\color{strings},
    basicstyle=\ttfamily\footnotesize,
    breakatwhitespace=false,         
    breaklines=true,                 
    captionpos=b,                    
    keepspaces=true,                 
    numbers=left,                    
    numbersep=5pt,                  
    showspaces=false,                
    showstringspaces=false,
    showtabs=false,                  
    tabsize=2
}
\lstset{style=python_style}

\title{Final Project}
\author{Luis Gilbuena, Matthew Merritt, Michael Merritt, Anton Petrenko}
\date{May 7, 2025}

\begin{document}

\maketitle

\section{Approach}

\subsection{Primary Data Structure}

The key idea to our approach was to construct a dictionary that would use $n$-grams (represented as tuples) as keys and return the most likely part of speech for the word given
that specific $n$-gram. This was done by going through the dataset and constructing $n$-grams from each sentence using the actual part of speech, and tallying the results as we
went. Once all of the dataset was parsed and added to the dictionary, the dictionary was then separated into two sets: 
\begin{itemize}
    \item \textbf{``Unanimous''} $n$-grams that always resulted in the same part of speech during training regardless of additional context, such as a bigram that always 
    resulted in the word being tagged as a verb.
    \item \textbf{``Non-unanimous''} $n$-grams that could result in different parts of speech with additional context, such a trigram that could result in being tagged as a noun
    or a verb based on the words before it.
\end{itemize}
Once these two sets are separated, we then store only the most likely option for the part of speech for every $n$-gram to return when queried, and remove all of the additional 
tallying information.

\subsection{$n$-gram Lookup Order}

Using the two sets mentioned above, we modified our approach to use the smallest amount of information needed to get a definitive answer.
Specifically, we first check the unanimous set to see if the unigram always leads to the same part of speech regardless of what comes next, followed by the same check for the 
bigram, trigram, and quadgram. Assuming that there was no single answer beforehand, we then use the non-unanimous set to find the most likely result from the quadgram, trigram, 
bigram, or unigram, checking them in that order. If there are still no results for the given $n$-gram in either set, we default to assuming that the word is a noun (``NN''), 
the most common part of speech.

\subsection{Jackknife Training}

For the training process, we used a jackknife technique to split the data into separate training and testing datasets, with each slice created by the jackknife procedure
setting aside 1000 sentences for the testing dataset, and using the remaining sentences as the training dataset. Since there were 47959 sentences in total, this means that
48 different training and testing sets were created and tested (the final testing set was only 959 sentences rather than the full 1000, since the dataset did not split evenly).
The accuracy was then taken as the average of all 48 different trials.

\subsection{Separate Implementations}
To test out a few different implementation ideas we had as a group, we created three separate implementations of the aforementioned approach.

\begin{itemize}
    \item \textbf{Luis's implementation}, which uses a JSON file to load the sentence data and uses sets to handle tiebreakers.
    \item \textbf{Anton's implementation}, which uses Python's built-in CSV module and strings as keys for the dictionaries instead of tuples.
    \item \textbf{The Merritt's implementation}, which uses pandas DataFrames to load the data all at once and then applies functions to columns and rows of the data to build our model.
\end{itemize}

Each implementation will be discussed more in depth in its own section below.\\

\section{Implementations}

This section focuses on the Python implementations of the general approach described above.

\subsection{Luis's Implementation}

\subsubsection{Implementation}

In cases where a word has multiple parts of speech with equal highest frequencies---say the word `accused` has a unigram frequency dictionary of \{`VBN`:3, `VBD`:3\}---we assign the word a set containing parts of speech tied for the highest frequency. In the case we reach the non-unanimous dictionaries, instead of checking if the true part of speech is equal to the value, given the key is in the dictionary, we check if the true part of speech is inside the set.

\begin{lstlisting}[language=Python]
for x in range(4):
        mdl = model[x]
        for key in list(mdl.keys()):
            if len(mdl[key]) > 1:
                maxvalue = max(mdl[key].values())
                maxval = set()
                for y in mdl[key]:
                    if mdl[key][y] == maxvalue:
                        maxval.add(y)
                multimodel[x][key] = maxval
                del mdl[key]
            else:
                pos = [x for x in model[x][key]]
                model[x][key] = pos[0]
    return model + multimodel
\end{lstlisting}

\subsubsection{Results}
The accuracy of this implementation of the $n$-gram tagger at best was 97\%, with total runtime of 378 seconds.

\subsection{Anton's Implementation}

\subsubsection{Implementation}

In Anton's implementation, there is only one dictionary for all the \textit{n}-grams.
When building the model, the unigram model is built first.
The reason for this is to be able to identify the unanimous words from the non-unanimous words.
Following the creation of the unigram model, another pass is done through the data.
In this pass the bigram, trigram, and quadgram models are created.
For efficiency, the second pass includes a lookup to the existing unigram model for each word.
If the word only ever results in one part of speech, there is no reason to save more data about this word.
If the word appears as multiple parts of speech, we construct each gram by inserting it into the existing dictionary with the history encoded in the key.
The resulting dictionary contains all \textit{n}-gram versions.
Keys in the dictionary are strings of the relevant words, where any relevant history is encoded before the word using a circumflex as a delimiter.


\begin{lstlisting}[language=Python]
for sentence in lines:
    history = []
    for word, pos in sentence:
        if len(unigram[word].keys) > 1:
            for i in range(len(history)):
                unigram[f"{'^'.join(history[-(i+1):])}^{word}"][pos] += 1
        history.append(pos)
        if len(history) == n: 
            history.pop(0)
\end{lstlisting}
\noindent This is a small pseudocode snippet denoting the process of building the \textit{n}-grams in the second pass of the data.
As shown, the unigram model is always checked before constructing the other \textit{n}-grams in the same dictionary.
Keys into the dictionary are strings which lead to another dictionary which contains the part of speech as keys and the number of times it appeared as that part of speech as the values.
Note the creation of keys which have the history encoded in them, as well as the clearing of the history for each sentence.
The array which holds the history is also bound by the highest number \textit{n} denoting the highest \textit{n}-gram model to create.

\subsubsection{Results}
The results of this implementation of the $n$-gram tagger were as follows:
\begin{itemize}
    \item \textbf{Unigram only:} 95.28\% accuracy, runtime of 1m 45.964s
    \item \textbf{Up to bigram:} 98.07\% accuracy, runtime of 1m 46.061s
    \item \textbf{Up to trigram:} 98.6\% accuracy, runtime of 1m 45.628s
    \item \textbf{Up to quadgram:} 98.93\% accuracy, runtime of 1m 45.106s
\end{itemize}

\subsection{The Merritt's Implementation}

\subsubsection{Implementation}

This approach utilizes the pandas Python library for data processing, reading in data as a tabular format and applying operations to entire columns at a time to achieve faster transformations. The implementation and can be divided into four parts: initial data processing, model generation, testing against a model, and jackknife training. Each part is implemented in a separate function to allow for flexibility in terms of $n$-gram size, excluding initial data processing which only occurs once and does not take any arguments. 

In in the initial data processing stage, the tagged sentence data is read from the spreadsheet, with each row containing a single word in a sentence and the expected part of speech, and with numbers marking the beginning of sentences. To simplify the process of dividing the dataset into sentences for training and testing, the sentence numbers are cascaded down to the subsequent words that occur before the next sentence starts. Additionally, the sentence number is extracted from the original format by stripping away non-numeric characters. Finally, string data types needed to be specified and missing value detection needed to be disabled when loading the dataset in using pandas to avoid problems with empty fields or words reserved for values in Python, i.e. None and False.

In the model generation stage, two major steps occur: mapping and flattening. Before mapping can occur, the history of previous parts of speech must be generated for each word and stored in the row for later use. To do this, the part of speech column is copied and shifted down based on the number of history entries required. For example, when generating quadgrams, the history requires the previous three parts of speech, meaning that three shifted columns need to be added. Values in the shifted columns are left blank if the sentence number changes after a shift. With the history added to each row in the table, a mapping can be generated by applying a function to the rows of the table, creating a dictionary that maps history patterns to occurrence counters for each part of speech. This step of applying the function to the rows to build the dictionary is the most time consuming step in our implementation, taking roughly 15 seconds per jackknife iteration. With the mapping dictionary created, a flattening step is performed to create the model where the mappings are separated into separate dictionaries for unanimous and non-unanimous n-grams that skip past the counts for each part of speech and instead return only the most probable part of speech. 

One modification explored but ultimately not implemented at this model generation stage was aimed at reducing ambiguity, but it was not included due to negative impact on accuracy. Specifically, in the case of the non-unanimous model, we explored removing cases where the part of speech decision was ambiguous. For example, when using the first thousand sentences as the testing set and the remaining sentences as the training set, one such ambiguity occurs for the the word 'ruling' after a preposition, proper noun, then possesive ending. In this case, there is an equal chance of 'ruling' representing a 'Verb, gerund or present participle' ('VBG') or a 'Noun, singular or mass' ('NN'), with both occurring 5 times in the training set. Removing this pattern from our model did not cause a noticeable increase in accuracy, so ambiguity is instead currently decided on a seemingly arbitrary basis. Currently, in the case of ties, the pattern that was observed first will be selected by the max function call on the dictionary of counts.

\begin{lstlisting}[language=Python]
    def generate_ngram_model(training_set, gram_size=4):
        # ---------------- STAGE 0: PREPROCESSING HISTORY ----------------
        modified = pd.DataFrame(training_set)
    
        for n in range(gram_size - 1):
            modified[f'Previous POS {n + 1}'] = modified['POS'].shift(n + 1)
            modified[f'Previous POS {n + 1}'] = modified[f'Previous POS {n + 1}'].where(modified['Sentence #'] == modified['Sentence #'].shift(n + 1))
    
        # ---------------- STAGE 1: INITIAL COUNTING ----------------
        ngram_map = {}
    
        def write_to_map(row):
            previous_pos = list(map(lambda n : row[f'Previous POS {n}'], list(range(1, gram_size))[::-1]))
            previous_pos = list(filter(lambda pos : not pd.isna(pos), previous_pos))
    
            patterns = [(*previous_pos[-i:], row['Word']) for i in range(1, len(previous_pos) + 1)]
            patterns.append(row['Word'])
    
            for pattern in patterns:
                if pattern in ngram_map:
                    if row['POS'] in ngram_map[pattern]:
                        ngram_map[pattern][row['POS']] += 1
                    else:
                        ngram_map[pattern][row['POS']] = 1
                else:
                    ngram_map[pattern] = { row['POS']: 1 }
    
        modified.apply(write_to_map, axis=1)
    
        # ---------------- STAGE 2: FLATTENING ----------------
    
        # index 0 is the unanimous map, index 1 is the highest probability map
        ngram_model = [{}, {}]
        for (key, pos_map) in ngram_map.items():
            map_index = 0 if len(pos_map) == 1 else 1
    
            ngram_model[map_index][key] = max(pos_map, key=pos_map.get)
    
        return (ngram_map, ngram_model)
\end{lstlisting}

In the testing against a model stage, the model from the previous stage is taken alongside the testing set, and each row is tested using both the unanimous and non-unanimous models when necessary. First the testing set is augmented with history similar to the training set, and then a function is applied on each row that tests and stores the outcomes into a new column. The testing order is to go through all the unanimous models from least amount of context (unigram) to most amount of context (quadgram in our case), then go through all the non-unanimous models in the opposite order. Output is appended to a text file for debugging and logging. This log file should be deleted between separate iterations of the program to prevent confusion.

In the final jackknifing stage, the raw dataset is divided into training and testing sets that can be passed to the model generation and testing stages. Each iteration selects 1000 sentences for testing and passes the remaining sentences as the training set. Between iterations, the entire model is discarded and rebuilt, with only the correct and incorrect result counts remaining across iterations. Overall accuracy is output at the end.

\subsubsection{Results}
Ultimately, this approach yielded lower than expected accuracy and generally performed slower than other approaches. In terms of the performance, the biggest impact was in the model generation stage when the mapping function is applied to the rows of the training set. We were able to iterate on this, having initially iterated over the rows individually, taking 2-3 minutes per training set. In the end, we reduced this down to about 15 seconds by applying a function to the rows of the table instead, but this still adds up. Regarding the lower accuracy, comparing our quadgram results to the provided sample output (99.00\% accuracy), we found that only additional incorrect tests we had were when ambiguous decisions occurred, making up the 4\% difference in accuracy. 

The results of this implementation of the $n$-gram tagger were as follows:
\begin{itemize}
    \item \textbf{Unigram only:} 93.60\% accuracy, runtime of 5m 13.980s
    \item \textbf{Up to bigram:} 95.84\% accuracy, runtime of 7m 46.116s
    \item \textbf{Up to trigram:} 95.94\% accuracy, runtime of 10m 46.101s
    \item \textbf{Up to quadgram:} 95.91\% accuracy, runtime of 13m 39.980s
\end{itemize}

\section{Reflection}

Some sort of conclusion goes here.

% \begin{center}
%     \begin{tabular}{ c | c | c }
%         \hline
%         & & Accuracy & Runtime (s)
%         \hline
%         anton & Unigram & 95.28\% & 17.720
%         & Bigram & 98.07\% & 26.394 
%     \end{tabular}
% \end{center}

\noindent
\begin{tabular}{|m{4cm}|m{4cm}|m{3cm}|m{3cm}|}
    \hline
    \textbf{} & \textbf{$n$-gram Size} & \textbf{Accuracy} & \textbf{Runtime (s)} \\
    \hline
    \multirow{4}{*}{Anton}  & Unigram & 95.28\% & 17.720 \\
                            \cline{2-4}
                            & Bigram & 98.07\% & 26.394 \\
                            \cline{2-4}
                            & Trigram & 98.60\% & 37.253 \\
                            \cline{2-4}
                            & Quadgram & 98.93\% & 48.513 \\
    \hline
    \multirow{4}{*}{Luis}   & Unigram & 94.45\% & 105.964 \\
                            \cline{2-4}
                            & Bigram & 96.67\% & 106.061 \\
                            \cline{2-4}
                            & Trigram & 96.86\% & 105.628 \\
                            \cline{2-4}
                            & Quadgram & 97.22\% & 105.106 \\
    \hline
    \multirow{4}{*}{Merritt}& Unigram & 93.60\% & 313.980 \\
                            \cline{2-4}
                            & Bigram & 95.84\% & 466.116 \\
                            \cline{2-4}
                            & Trigram & 95.94\% & 646.101 \\
                            \cline{2-4}
                            & Quadgram & 95.91\% & 819.980 \\
    \hline
\end{tabular}

\end{document}
